\documentclass[10pt,a4paper]{amsart}
\usepackage[margin=19mm]{geometry}
\usepackage{amsmath,amssymb,mathtools}
\usepackage[scr=rsfs,cal=euler]{mathalfa}
\usepackage{xfrac}
\usepackage[unicode,hidelinks,bookmarks=false]{hyperref}
\newcommand{\ie}{i.e.\ }
\newcommand{\cf}{cf.\ }
\newcommand{\resp}{respectively\ }
\newcommand{\Subsection}{\subsection}
\providecommand{\nobreakdash}{\nobreak\hbox{-}}
\setlength{\emergencystretch}{2em}
\multlinegap0pt
\usepackage{bm}
\usepackage{bbm}
\usepackage{dsfont}
\usepackage{xspace}
\usepackage{cancel}
\usepackage{mathtools}
\usepackage{amsthm}
\makeatletter
\@ifundefined{lemma}{\newtheorem{lemma}{Lemma}}{}
\makeatother
\newcommand{\PP}{\mathbb{P}}
\newcommand{\QQ}{\mathbb{Q}}
\title[Rational points on K3 surfaces of degree 2]{Rational points on K3 surfaces of degree 2}
\author{J\'ulia Mart\'inez-Mar\'in}
\date{Source: arXiv:2505.13262v2 (CC BY 4.0)}
\begin{document}
\maketitle

\noindent\emph{Source and licence.} Rational points on K3 surfaces of degree 2. Version arXiv:2505.13262v2 (CC BY 4.0), distributed under the Creative Commons Attribution 4.0 International licence, as recorded in the arXiv licence metadata for this exact version and shown on the abstract page https://arxiv.org/abs/2505.13262v2. Full rights notice: licenses/arxiv-2505.13262-CC-BY-4.0.txt.\par\smallskip

\noindent\textit{Editorial scope.} Counted unit: the Lemma whose source label is \texttt{dominant} in the article (source0.tex lines 362--364, source label \texttt{dominant}), delivered together with the article's own complete original proof of it (source0.tex lines 365--383, one \texttt{proof} environment of 19 lines). No mathematics is written, repaired or re-derived: statement and proof are the article's own text line for line. Required context retained uncounted: none. Every definition, display and label of those lines is retained unchanged. Omitted from this package and not redistributed: 1--361, 384--477. Nothing inside a retained excerpt is omitted or altered: every retained body is delivered exactly as the article's lines are written, so no omission receipt of any kind accompanies this package. Citations: the retained excerpts carry no citation command at all, so no bibliography entry of the article is transcribed and no reference is redistributed. Reference adaptation: none was needed and none was made; every reference inside the delivered unit resolves inside it, so no unresolved reference or citation is produced. Printed slips kept literal: no printed word, symbol, hypothesis or number of the counted statement or of its proof was altered. Layout and class: the article's own class is \texttt{amsart} with packages including geometry, hyperref, graphicx, framed, amsmath, amssymb, amsfonts, mathrsfs, array, enumerate, enumitem, mathtools, tikz-cd, multirow; the wrapper is the lane's pinned \texttt{amsart} editorial wrapper at 10pt on A4, which restates the theorem-like environments the retained text needs and adds the article's own macro definitions that the retained text needs (\texttt{\textbackslash PP}, \texttt{\textbackslash QQ}), with extra wrapper packages (bm, bbm, dsfont, xspace, cancel, mathtools). The delivered numbering is the unit's own. Assets: no retained span contains a figure environment, a graphics-inclusion command, a PDF-page inclusion, a table or a TikZ picture; no image file is copied into this package, the package declares no asset dependency and no image file is redistributed with it. Licence: version arXiv:2505.13262v2 of this article is distributed under the Creative Commons Attribution 4.0 International licence, as recorded in the arXiv licence metadata for this exact version and shown on the abstract page https://arxiv.org/abs/2505.13262v2. A full rights notice is delivered at licenses/arxiv-2505.13262-CC-BY-4.0.txt. Characters: every retained excerpt is pure ASCII, so the delivered engine, XeLaTeX with its default Latin Modern text font, reports no missing character.
\begin{lemma} \label{dominant}
The restriction map $\psi|_{W}$ is dominant. 
\end{lemma}

\begin{proof}
It suffices to check this over the algebraic closure $\overline{\QQ}$. Let $\mathcal{U}\subset \mathcal{M}_2$ be the open dense subset with the property above. Let $X\in \psi^{-1}(\mathcal{U})$ be a K3 surface of degree 2 and let $X\rightarrow \PP^2$ be a double cover branched along a smooth sextic curve $B$. Let $P\in B(\overline{\QQ}) $ and let $\ell$ denote the tangent line to $B$ at $P$. Then there exists a linear transformation $A_1$ sending $P$ to the point $[1:0:0]$ and $\ell$ to the line $V(y)\subset \PP^2$, so $[1:0:0:0]\in A_1(X)(\overline{\QQ})$ and $V(y)$ is tangent to the branch curve of $A_1(X)$ at $[1:0:0]$. This implies that, in the equation defining the K3 surface $A_1(X)$, the coefficients of $x^6$ and $x^5z$ vanish and we can assume that the coefficient of $x^5y$ is $1$. Let $\beta_1, \beta_2, \beta_3 \in \overline{\QQ}$ denote the coefficients of $x^4y^2$, $x^4yz$ and $x^4z^2$ in the equation defining $A_1(X)$, respectively. We claim that there exists a linear transformation $A_2$ sending $A_1(X)$ to a surface such that its defining equation has $0$, $11$, $7$, $1$, $5$ and $7$ as the coefficients of $x^6$, $x^5y$, $x^5z$, $x^4y^2$, $x^4yz$ and $x^4z^2$, respectively.
Once we prove this claim, we are done: we can take the linear transformation $A_3$ sending the variables $x$, $y$ and $z$ to $\sqrt[6]{7/73}x$, $\sqrt[6]{7/73}y$ and $\sqrt[6]{7/73}z$, respectively, and we have that $Y \coloneqq A_3(A_2(A_1(X)))\in V(\overline{\QQ})\cap M'(\overline{\QQ})=W(\overline{\QQ})$, so $\psi|_{W}(Y)=\psi(X)\in \mathcal{U}$. 

To prove the claim, consider a transformation $A_2$ given by
\begin{align*}
x \mapsto &  \ x+ay+bz, \\
y \mapsto & \  11y+7z, \\
z \mapsto & \ cy+dz,
\end{align*}
for some $a,b,c,d\in \overline{\QQ}$. We want to find $a,b,c,d\in \overline{\QQ}$ such that the corresponding transformation $A_2$ is as claimed. Note that, to look at the coefficients of $x^6$, $x^5y$, $x^5z$, $x^4y^2$, $x^4yz$ and $x^4z^2$ in the equation defining the surface $A_2(A_1(X))$, we only need to check how $A_2$ acts on the coefficients of those same six monomials in the equation defining $A_1(X)$, by our choice of transformation $A_2$. Since the coefficient of $x^6$ in the equation of $A_1(X)$ is zero, the coefficient of $x^6$ in the equation of $A_2(A_1(X))$ also vanishes. Looking at how the other coefficients transform under the action of $A_2$, we have that $A_2(A_1(X))$ is as claimed if the following three equations hold.

\begin{align} 
    1 = & \ 121\beta_1+11\beta_2c+\beta_3c^2+55a, \label{x^4y^2} \\
    5 = & \ 154\beta_1+7\beta_2c+11\beta_2d+2\beta_3cd+35a+55b, \label{x^4yz} \\
    7 = & \ 49\beta_1+7\beta_2d+\beta_3d^2+35b. \label{x^4z^2}
\end{align}
Isolating $a$ and $b$ in equations \eqref{x^4y^2} and \eqref{x^4z^2} and substituting them in equation \eqref{x^4yz} yields a polynomial equation in the two variables $c,d$. Set, for example, $c=1$ and choose $d\in \overline{\QQ}$ to be a root of the resulting polynomial in one variable. This gives a transformation $A_2\in \operatorname{GL}(3, \overline{\QQ})$ such that $A_2(A_1(X))$ is as claimed, and this finishes the proof. 
\end{proof}

\end{document}
